\documentclass{tikzposter}
\usepackage[utf8]{inputenc}
\usepackage[T1]{fontenc}
\usepackage[french]{babel}
\usepackage{amsmath,amssymb}
\usepackage{graphicx}
\usepackage{listings}
\usepackage{xcolor}

\lstset{
    language=Python,
    basicstyle=\ttfamily\footnotesize,
    keywordstyle=\color{blue},
    commentstyle=\color{green},
    stringstyle=\color{red},
    numbers=left,
    numberstyle=\tiny,
    frame=single,
    breaklines=true
}

\title{Architecture Technique: Motion Flow pour Animation d'Images}
\author{Arthur - Détails d'Implémentation}
\institute{Makeitalive - Generative AI Project}

\begin{document}
\maketitle

\begin{columns}
\begin{column}{0.5}
\block{Architecture du Modèle MotionFlowUNet}{
Le modèle principal est basé sur une architecture U-Net modifiée pour la prédiction de champs optiques :

\textbf{Structure:}
\begin{itemize}
\item \textbf{Entrée:} Image RGB 3 canaux (512x512)
\item \textbf{Encodeur:} 4 niveaux de convolution + MaxPool
\item \textbf{Bottleneck:} 1024 canaux après downsampling x16
\item \textbf{Décodeur:} 4 niveaux d'upsampling + concaténation
\item \textbf{Sortie:} Champ de mouvement 2 canaux (dx, dy)
\end{itemize}

\textbf{Détails techniques:}
\begin{itemize}
\item DoubleConv: Conv(3x3) + BN + ReLU répété
\item Down: MaxPool2d + DoubleConv
\item Up: Bilinear Upsample + Conv + Concat skip
\item Pas de BatchNorm dans les dernières couches
\end{itemize}
}

\block{Implémentation en PyTorch}{
\lstinputlisting[language=Python, firstline=25, lastline=70]{../src/motion_flow/model.py}
}

\block{Processus d'Entraînement}{
\textbf{Configuration:}
\begin{itemize}
\item Optimizer: Adam (lr=1e-4)
\item Loss: MSE entre flow prédit et ground truth
\item Batch size: 8
\item Époques: 100
\item Device: CUDA si disponible
\end{itemize}

\textbf{Dataset LandscapeMotionDataset:}
\begin{itemize}
\item Paires d'images consécutives
\item Calcul du flow ground truth avec OpenCV
\item Augmentations: resize, normalize
\item DataLoader avec num\_workers=16
\end{itemize}
}
\end{column}

\begin{column}{0.5}
\block{Génération d'Animation}{
\textbf{Algorithme d'inférence:}
\begin{enumerate}
\item Charger image statique
\item Prédire flow avec le modèle
\item Appliquer warping avec scipy.ndimage.map\_coordinates
\item Générer séquence d'images animées
\item Créer vidéo/GIF avec OpenCV
\end{enumerate}

\textbf{Code d'application du flow:}
\begin{lstlisting}[language=Python]
# Application du motion flow
map_y = y - flow_y
map_x = x - flow_x

warped = np.zeros_like(image)
for channel in range(3):
    warped[..., channel] = map_coordinates(
        image[..., channel],
        [map_y, map_x],
        order=1, mode='reflect'
    )
\end{lstlisting}
}

\block{Comparaison des Approches}{
\textbf{LoRa Fine-tuning vs From Scratch:}

\begin{tabular}{|l|c|c|}
\hline
Aspect & LoRa & Motion Flow \\
\hline
Complexité & Faible & Moyenne \\
\hline
Données requises & Beaucoup & Peu \\
\hline
Contrôle & Limité & Complet \\
\hline
Qualité & Variable & Consistante \\
\hline
Temps d'entraînement & Court & Long \\
\hline
\end{tabular}

\textbf{Avantages Motion Flow:}
\begin{itemize}
\item Entraînement from scratch plus léger
\item Contrôle direct du mouvement
\item Pas de dépendance à des modèles pré-entraînés
\item Architecture optimisée pour la tâche
\end{itemize}
}

\block{Métriques et Évaluation}{
\textbf{Loss Function:}
\[ \mathcal{L} = \frac{1}{N} \sum_{i=1}^{N} \| \hat{f}_i - f_i \|^2 \]
Où $\hat{f}_i$ est le flow prédit et $f_i$ le ground truth.

\textbf{Visualisation:}
\begin{itemize}
\item Champs de mouvement colorés (coolwarm)
\item Comparaison avant/après warping
\item Génération de GIFs animés
\item Analyse des erreurs par région
\end{itemize}
}

\block{Défis et Solutions}{
\textbf{Problèmes rencontrés:}
\begin{itemize}
\item Artefacts aux bords (solution: mode='reflect')
\item Mouvement irréaliste (solution: régularisation)
\item Overfitting (solution: data augmentation)
\end{itemize}

\textbf{Optimisations futures:}
\begin{itemize}
\item Architecture plus profonde
\item Loss functions avancées (SSIM, perceptual)
\item Conditionnement temporel
\item Multi-resolution training
\end{itemize}
}
\end{column}
\end{columns}

\end{document}